\section{问题分析}

\subsection{问题一分析}

问题一的核心矛盾并非单纯“最大程度去噪”，而是“在抑制伪影的同时尽可能保留未知的视觉任务特征”。若直接采用强平滑、固定模板回归或过度分解，虽然可降低信号方差，却可能同时削弱P300等晚期成分以及左右刺激之间本就微弱的差异。

因此本文将问题一拆分为四个递进环节：

\begin{enumerate}[label=(\arabic*)]
    \item 建立低偏差的基础预处理，获得可比较的保真参考；
    \item 根据试次质量构造稳健权重，降低异常试次对平均响应的影响；
    \item 构造受约束的候选降噪算子，并通过已知信号注入实验检验“去噪--保真”权衡；
    \item 在冻结的主处理方案上提取ERP，分别分析共同视觉响应与左右刺激差分，并进行低维曲线拟合。
\end{enumerate}

\subsection{问题二分析}

问题二要求从“观测差异”进一步过渡到“生成机制”。因此不能仅依赖分类器，而需要显式建立LGN输入、皮层神经群体状态以及头皮电极观测之间的映射。后续模型应同时满足：具有明确状态变量、能够产生与实测EEG可比较的时间序列、能够解释左右视觉刺激在空间配置或输入结构上的差异，并可通过留出数据验证。

\subsection{问题三分析}

\TODO{根据问题三完整题意与后续代码补充。建议强调：从局部视觉响应机制进一步上升到认知/记忆状态的动态建模，并与问题一、二共享观测接口。}
